% ============================================================================
%  Section 4: Measuring the Value of Retirement Reforms
%  Sources: FROT 2025_sub/12_app_quant.tex ("Construction of Reforms",
%  "Construction of Outcomes"), 2025_sub/7_policy_reforms.tex; blueprint sec. 7.
%  Derivations: Appendix E (sections/E_derivations.tex).
% ============================================================================
\section{Measuring the Value of Retirement Reforms}\label{sec:framework}

The marginal value of public funds (MVPF) of a policy is the willingness to pay of its beneficiaries divided by its net cost to the government, including the fiscal externalities of the behavioral responses it induces \citep{hendren2016policy,hendren2020unified}. For a life-cycle policy, both objects must be measured over the whole life cycle: the willingness to pay before individuals know the health and earnings histories that determine what they receive, and the net cost across every program the policy affects, including responses that occur decades before or after the reformed rule applies. This section defines both objects in our model, relates their ratio to the sufficient-statistics analysis of pension reform, and derives the accounting identity that governs the fiscal effects of claiming-age reforms.

\subsection{Reforms}

A reform is a permanent change in the policy rules of Section~\ref{sec:model}, announced before the cohort enters the model at age 30 and in force over its entire life; the status quo is the estimated economy of Section~\ref{sec:estimation}. For each reform, we solve the individual's problem \eqref{eq:value} under the new rules and simulate the cohort from the same initial distribution. Because the reform is known from age 30, its effects include responses long before the reformed rule applies: a worker who expects a more generous benefit formula adjusts her hours at 35, when her earnings enter her AIME, as well as her exit at 63.

We study 66 reforms in five families, constructed in Appendix~\ref{app:reforms}. The first family (17 reforms) changes the claiming-age adjustment $\alpha(\cdot)$ of Section~\ref{sec:model_ss}. Reforms AN1--AN3 raise the NRA from 65 to 66, 67, and 68, which shifts the entire schedule of benefits by claiming age to the right. Reforms AD1--AD5 raise the DRC from 4.5 percent to 8, 8.5, 9, 9.5, and 10 percent and leave the reductions for early claiming unchanged. Reforms AL1--AL9 make the adjustment symmetric around the NRA of 65: the benefit changes by $x$ percent of the PIA for each year by which the claim precedes or follows the NRA, with $x$ from 7 to 15 percent. \autoref{fig:schedules} and \autoref{tab:schedules} display the schedules. The second family (39 reforms) raises the replacement rates of the benefit formula \eqref{eq:pia}. Reforms P1.1--P1.9 raise the first rate from 0.90 to 0.91, \ldots, 0.99; with the first rate at 0.99, reforms P2.0--P2.9 raise the second rate from 0.32 in ten steps of 1 percent of its initial value, to 0.3232, \ldots, 0.352; and with the second rate at 0.352, reforms P3.0--P3.9 raise the third rate from 0.15 in the same proportional steps, to 0.1515, \ldots, 0.165. Reforms L0--L9 raise all three rates in proportion, by 1 to 10 percent.

The third family (two reforms) raises the share of medical expenditures after age 65 that Medicare covers by 20 percentage points, for all individuals (Part A) or for labor market nonparticipants (Part B), relative to a baseline in which coverage rates vary with health and working status as in the MEPS. The fourth family (five reforms) lowers income taxes in old age: reforms I1--I5 raise the level parameter of the tax function \eqref{eq:tax} from $\lambda_t$ to $\lambda_t+x(t-61)$ at ages $t\geq 62$, with $x$ from 1 to 5 percent, which raises after-tax income at every income level by the proportion $x(t-61)/\lambda_t$. The fifth family (three reforms) raises the net return on assets at every age by 0.2, 0.4, and 0.8 percentage points; the largest subsidy sets the net return equal to the pretax return $r_f$. All reforms except the savings subsidies change benefits, coverage, or taxes only from age 62, the earliest claiming age; this distinction governs the conversion of welfare into dollars (Section~\ref{sec:fw_mvpf}).

We also remove the existing instruments one at a time: D1 sets Social Security benefits to zero, D2 compares the status quo with full coverage of medical expenditures after age 65, D3 sets the net return equal to the pretax return, D4 eliminates the income tax, and D5 eliminates DI. These experiments are not financed, and we report only their welfare effects (Section~\ref{sec:results_elim}).

\subsection{Welfare}\label{sec:fw_welfare}

Let $\bar V^{0}$ be the average over simulated individuals of the value at age 30, $V_{30}(X_{30})$, in the status quo, and $\bar V^{1}$ the same average under a reform. Flow utility \eqref{eq:utility} is homogeneous of degree $\gamma(1-\nu)$ in consumption: multiplying consumption by $e^{x}$ at every age and in every state, with leisure unchanged, multiplies flow utility by $e^{\gamma(1-\nu)x}$, and therefore multiplies each individual's expected lifetime utility, and their average $\bar V^{0}$, by the same factor. The consumption-equivalent welfare change (CE) of a reform is the $x$ that solves $\bar V^{1}=e^{\gamma(1-\nu)x}\,\bar V^{0}$:
\begin{equation}\label{eq:ce}
\mathit{CE}=\frac{\ln\big(\bar V^{1}/\bar V^{0}\big)}{\gamma(1-\nu)},
\qquad
\mathit{LE}=\frac{\ln\big(\bar V^{1}/\bar V^{0}\big)}{(1-\gamma)(1-\nu)}.
\end{equation}
The ratio $\bar V^{1}/\bar V^{0}$ is positive because, with $\nu>1$, lifetime utility is negative under every policy. $\mathit{CE}$ is the proportional increase in status-quo consumption, at every age and in every state, that the cohort values as much as the reform. The leisure equivalent $\mathit{LE}$, obtained by scaling leisure instead, equals $\mathit{CE}\,\gamma/(1-\gamma)$ and ranks reforms identically. We report $\mathit{CE}$, in percent, as the willingness to pay for the reform, $\WTP$. As a log ratio, it is antisymmetric: the CE of a move from one set of rules to another is minus the CE of the reverse move.

The willingness to pay is ex ante. It is evaluated at age 30, before health, ability, and survival histories are realized, and it averages expected utility with equal weights across individuals who differ in initial assets, wages, and health. It therefore includes the value of a reform as insurance. The welfare effect of a change in resources in a given state is weighted by the marginal utility of consumption in that state, so a reform that raises resources in states where marginal utility is high is worth more ex ante than its expected transfer, and a reform that lowers them in such states costs more than its expected saving. Health shocks, low ability, and medical expenses determine which states these are; the claiming-age schedule and the benefit formula determine the resources available in them.

Many applications of the MVPF instead sum the ex post willingness to pay of the individuals a policy affects, each in her own dollars, which for small policies equals the mechanical transfer each receives; distributional and insurance considerations then enter through separate welfare weights \citep{hendren2020unified,Finkelstein2020welfare,saez2016generalized}. Our measure embeds the weights of a utilitarian evaluation of one cohort behind the veil of ignorance. It does not weigh the welfare of different cohorts (Section~\ref{sec:financing}).

\subsection{Fiscal Accounting}\label{sec:fw_fiscal}

The fiscal effects of a reform are changes in the government accounts of Section~\ref{sec:model_gov}: the per capita present values at age 62, discounted at the net return $r$, of Social Security retirement outlays $\mathit{SS}$, DI outlays $\mathit{DI}$, consumption-floor transfers $\mathit{TR}$, income-tax revenue $\mathit{IT}$, payroll-tax revenue $\mathit{PT}$, and Medicare payments, with net government spending $G=\mathit{SS}+\mathit{DI}+\mathit{TR}-\mathit{IT}-\mathit{PT}$. Flows before age 62, such as the taxes paid during the working life, are compounded forward to 62. We denote by $\Delta$ the difference between the reform and the status quo. No reform is required to balance the budget: the MVPF values the change in net cost without specifying how it is financed or how a saving is used.

We split the net cost of a reform into a direct cost on the program that the reform changes and an indirect cost on all other programs. The direct cost includes the behavioral responses within the reformed program: for a claiming-age reform, it is the change in retirement outlays, which combines the change in benefits at given claiming ages with the effect of the induced changes in claiming ages. The indirect cost collects the fiscal externalities on all other programs: income and payroll taxes, DI, transfers, and, for reforms outside Social Security, Social Security itself. In dollars per capita, the direct cost $\mathit{DC}$, the indirect cost $\mathit{IC}$, and the net cost $\mathit{NC}=\mathit{DC}+\mathit{IC}$ are
\begin{alignat*}{4}
&\text{Social Security:}\quad &&\mathit{DC}=\Delta\mathit{SS}, \quad &&\mathit{IC}=\Delta(\mathit{DI}+\mathit{TR}-\mathit{IT}-\mathit{PT}), \quad &&\mathit{NC}=\Delta G,\\
&\text{Medicare:} &&\mathit{DC}=\Delta\mathit{MC}^{tg}, &&\mathit{IC}=\Delta G, &&\mathit{NC}=\Delta G+\Delta\mathit{MC}^{tg},\\
&\text{Income tax:} &&\mathit{DC}=-\Delta\mathit{IT}, &&\mathit{IC}=\Delta(\mathit{SS}+\mathit{DI}+\mathit{TR}-\mathit{PT}), \quad &&\mathit{NC}=\Delta G,\\
&\text{Savings subsidy:} &&\mathit{DC}=K, &&\mathit{IC}=\Delta G, &&\mathit{NC}=\Delta G+K,
\end{alignat*}
where the Social Security family comprises the claiming-age, progressivity, and level reforms, $\mathit{MC}^{tg}$ is the present value of Medicare's payments to the beneficiaries that the reform targets (all individuals aged 65 and older for Part A, nonparticipants aged 65 and older for Part B), and $K$ is the present value of the savings subsidy, the additional return the government pays on assets from age 30.

% AUTHORS: the absolute-value bases are inferred from the tables, not stated in FROT 2025_sub/12_app_quant.tex. Regressing Net on Direct and Indirect fits to rounding precision only with |G0| as base and G0<0 (SS0 ~ 0.137|G0|, IT0 ~ 1.13|G0|, PT0 ~ 0.076|G0| in the Medicare baseline). Confirm against the code.
The results tables report each cost as a percentage of the absolute value of a baseline present value, so that a positive entry is a cost to the government and a negative entry a saving. Net is expressed relative to baseline net government spending, $|G^{0}|$. Direct is expressed relative to baseline retirement outlays $\mathit{SS}^{0}$ for Social Security reforms, baseline income-tax revenue $\mathit{IT}^{0}$ for income-tax reforms, baseline payroll-tax revenue $\mathit{PT}^{0}$ for Medicare reforms---so that it is the proportional increase in the payroll tax that would finance the additional Medicare payments---and $|G^{0}|$ for savings subsidies. Indirect is expressed relative to the baseline value of the aggregate it measures (Appendix~\ref{app:outcomes}). Because the bases differ, Net is in general not the sum of Direct and Indirect; the dollar amounts add up.

\subsection{The Marginal Value of Public Funds}\label{sec:fw_mvpf}

We convert the willingness to pay into dollars with a present value of baseline per capita consumption $\bar c^{\,0}_{t}$:
\begin{equation}\label{eq:mvpf}
\MVPF=\frac{\WTP^{\$}}{\mathit{NC}},
\qquad
\WTP^{\$}=\frac{\WTP}{100}\,\bar C^{0}_{s},
\qquad
\bar C^{0}_{s}=\sum_{t=s}^{T}(1+r)^{-(t-62)}\,\bar c^{\,0}_{t},
\end{equation}
with $s=62$ for the reforms of Social Security, Medicare, and income taxes and $s=30$ for the savings subsidies. Numerator and denominator are per capita present values at age 62, in dollars.

We follow the conventions of \citet{hendren2020unified}. When the net cost is positive, the MVPF is the welfare delivered per dollar of government spending. When the WTP is positive and the net cost is zero or negative, the MVPF is infinite: the reform pays for itself, since the government could implement it, return the saving to the cohort, and keep the entire welfare gain. When the WTP and the net cost are both negative---a cut that lowers welfare and saves money---the MVPF is the welfare cost per dollar of revenue raised, the shadow price of raising revenue from the beneficiaries of a policy by reducing spending on it. By the antisymmetry of the CE and of the dollar fiscal effects, it equals the MVPF of the reverse reform up to the baseline used for the dollar conversion. Financing a reform with MVPF $\MVPF_1$ by a cut with MVPF $\MVPF_2$ raises welfare when $\MVPF_1>\MVPF_2$, the comparison Section~\ref{sec:financing} uses to order the ways of closing a financing gap.

For some reforms, we also report a lower bound that counts only the direct cost:
\begin{equation}\label{eq:mvpf_lb}
\underline{\MVPF}=\frac{\WTP^{\$}}{\mathit{DC}}.
\end{equation}
For a reform with positive WTP, negative net cost, and positive direct cost, $\underline{\MVPF}$ is the MVPF that would obtain if none of the fiscal externalities on other programs, which are savings, materialized; it bounds the infinite MVPF from below and isolates the part of the evaluation that rests on the reformed program alone. When the direct cost is negative as well, the MVPF is infinite under either cost. For the savings subsidies, the model taxes interest income as it accrues, as if all capital income were realized and taxed every period, whereas saving for retirement, much of it in retirement accounts, is long-term investment whose returns are taxed later or not at all. % Source: FROT 2025_sub/7_policy_reforms.tex (savings subsidy discussion: realization assumption, IRAs and 401(k) plans).
The indirect saving is therefore an upper bound, and we report the MVPF of a subsidy as the interval $[\WTP^{\$}/\mathit{DC},\,\WTP^{\$}/\mathit{NC}]$.

The dollar conversion is a convention. The CE scales consumption at every age from 30 to $T$, whereas \eqref{eq:mvpf} values it over the ages at which the reform operates: from 62 for reforms that start at 62, and over the whole life cycle for savings subsidies. It does not affect whether an MVPF is finite or infinite, which depends only on the signs of WTP and net cost, nor the ranking of reforms within a family, which share the base. It does affect the levels of finite MVPFs, which scale with the base: valuing the CE of a reform that starts at 62 over the whole life cycle would multiply its finite MVPF by $\bar C^{0}_{30}/\bar C^{0}_{62}$, and comparisons of MVPF levels with one, the MVPF of a transfer that induces no behavioral response, depend on the base in the same way. Because the savings subsidies use the larger lifetime base, the levels of their MVPFs are not directly comparable with those of the other families; Section~\ref{sec:discussion} quantifies the difference.

The construction implies a simple relation between the columns of the results tables. Since $\WTP^{\$}$ and $\mathit{NC}$ scale fixed baseline aggregates and $\mathit{Net}=100\,\mathit{NC}/|G^{0}|$,
\begin{equation*}
\MVPF=\frac{\WTP}{\mathit{Net}}\times\frac{\bar C^{0}_{s}}{|G^{0}|},
\end{equation*}
where the second factor depends on the baseline but not on the reform. Computed from the rounded table entries, this factor lies between 0.38 and 0.40 for all 48 reforms with a finite MVPF that start at 62 and between 1.87 and 1.99 for the three savings subsidies, a dispersion of the order of the rounding error in the entries. % Derived: MVPF x Net / WTP from tables/tab_claiming_body.tex, tab_generosity_body.tex, tab_generosity_app_body.tex, tab_other_body.tex (savings: upper end of interval).
Every MVPF computed with the net cost can thus be recovered, up to rounding, from the WTP and Net columns.

\subsection{Decomposition and Sufficient Statistics}\label{sec:fw_decomp}

The net cost of a reform also splits into a mechanical and a behavioral component. Let $\mathit{ME}$ be the change in net cost that the reform would cause if every individual kept her status-quo hours, claiming age, and saving, and let the fiscal externality $\mathit{FE}$ be the remainder, due to the behavioral responses. Then
\begin{equation}\label{eq:netcost}
\mathit{NC}=\mathit{ME}+\mathit{FE}=\mathit{DC}+\mathit{IC}.
\end{equation}
The direct cost contains the mechanical cost on the reformed program together with the behavioral responses within it; the indirect cost contains the fiscal externalities on all other programs together with mechanical interactions between programs, such as the fall in consumption-floor transfers when benefits rise for individuals at the floor (Appendix~\ref{app:derivations}).

For a small reform, the envelope theorem implies that behavioral responses have no first-order effect on welfare: since individuals were optimizing, the WTP equals the mechanical change in each individual's resources weighted by the marginal utility of consumption in the state in which it occurs \citep{hendren2016policy}. In our units, the CE of a small reform is the ratio of this marginal-utility-weighted transfer to the marginal-utility-weighted consumption stream of the status quo (Appendix~\ref{app:derivations}); the result holds with discrete participation and claiming choices because the value function is the upper envelope of the values of the alternatives \citep{milgrom2002envelope}. Behavioral responses then affect the MVPF only through the budget. We do not rely on this approximation: we compute WTP exactly from the value functions and $\mathit{FE}$ from the simulated responses.

The fiscal externality has three sources. The first is participation and hours: a worker who enters the labor force or postpones her exit pays income and payroll taxes, may forgo DI and transfers through their income and means tests, may have benefits withheld by the earnings test, and may raise her AIME and future benefits. Each such response lowers the net cost by the participation tax on her earnings: the taxes paid, benefits withheld, and DI and transfers forgone, net of the increase in future benefits \citep{kleven2006marginal}. The second is the claiming age, whose effect on retirement outlays Section~\ref{sec:fw_claiming} characterizes. The third is saving, which raises taxable interest income and lowers means-tested transfers. The fiscal columns are organized by program rather than by margin, but the margins map into programs: for a Social Security reform, claiming responses enter Direct, and participation and saving responses enter Indirect, except for their effects on retirement outlays through the earnings test and the AIME.

\citet{haller2022welfare} evaluates pension reforms with a dynamic version of the Baily--Chetty formula, which compares the value of the consumption smoothing a reform provides, per dollar of mechanical transfer, with a fiscal multiplier, the ratio of the reform's total fiscal cost to its mechanical cost; a reform is desirable when the former exceeds the latter. With the multiplier $\mathit{FM}\equiv\mathit{NC}/\mathit{ME}$,
\begin{equation*}
\MVPF=\frac{\WTP^{\$}/\mathit{ME}}{\mathit{FM}}:
\end{equation*}
the MVPF is the value of consumption smoothing per mechanical dollar divided by the fiscal multiplier. A reform is desirable when its MVPF exceeds that of the instrument that finances it, and the Baily--Chetty comparison is the case in which that instrument has an MVPF of one (Appendix~\ref{app:derivations}). For a reform with a positive mechanical cost, an infinite MVPF corresponds to a negative multiplier: the behavioral responses more than finance the mechanical transfer.

The structural model measures the fiscal externality over the whole life cycle, across every program, and by health, and it computes the ex ante WTP of reforms outside the range of historical variation. Quasi-experimental evidence provides credible local responses without a full model: of retirement to financial incentives \citep{Liebman2009labor,Manoli2016nonparametric}, of employment to the earnings test \citep{gelber2020employment}, of disability enrollment to the NRA \citep{Duggan2007aching}, and of retirement to statutory ages \citep{seibold2021reference}. Such responses can discipline and test those of a structural model, which in turn extends them to the life cycle and to every program; the two approaches are complements \citep{Chetty2009sufficient,Kleven2021sufficient}.

\subsection{The Fiscal Logic of Claiming-Age Reforms}\label{sec:fw_claiming}

The fiscal effects of claiming-age reforms rest on an accounting identity. Consider an individual alive at age 62, with a fixed primary insurance amount $\mathit{PIA}$, who claims at age $a$ at or after the NRA and receives $b(a)=\mathit{PIA}\,[1+\delta(a-\mathit{NRA})]$ every year from $a$ until death, where $\delta$ is the DRC. We abstract from the earnings test. Let $\mathcal{S}(t)$ be the probability, as of age 62, that she survives to age $t$, with $\mathcal{S}(62)=1$, and let the government discount at the net return $r$.

\begin{lemma}\label{lem:claiming}
Let $A(a)\equiv\sum_{t=a}^{T}(1+r)^{-(t-62)}\,\mathcal{S}(t)$. The present value at 62 of the individual's benefits when she claims at age $a\in\{\mathit{NRA},\ldots,69\}$ is $B(a)=b(a)\,A(a)$, and
\begin{equation}\label{eq:claimnpv}
B(a+1)-B(a)=\delta\,\mathit{PIA}\,A(a+1)\;-\;b(a)\,(1+r)^{-(a-62)}\,\mathcal{S}(a).
\end{equation}
A one-year delay therefore lowers the present value of her benefits if and only if $\delta<\delta^{*}(a)$, where
\begin{equation}\label{eq:faircredit}
\delta^{*}(a)\equiv\frac{b(a)\,(1+r)^{-(a-62)}\,\mathcal{S}(a)}{\mathit{PIA}\,A(a+1)}
=\frac{b(a)/\mathit{PIA}}{\sum_{t=a+1}^{T}(1+r)^{-(t-a)}\,\mathcal{S}(t)/\mathcal{S}(a)}.
\end{equation}
The rate $\delta^{*}(a)$ is strictly decreasing in the conditional survival probability $\mathcal{S}(t)/\mathcal{S}(a)$ at every age $t>a$ and nonincreasing in the terminal age $T$.
\end{lemma}

\noindent Appendix~\ref{app:derivations} gives the proof. The first term of \eqref{eq:claimnpv} is the cost of the higher benefit, paid for the rest of the individual's life; the second is the saving on the year of benefits she forgoes. The rate $\delta^{*}(a)$ is the credit that is actuarially fair from the government's perspective for this individual: it makes the program indifferent to her delaying by one year, given the government's discount rate and her survival prospects. At the NRA, where $b(a)=\mathit{PIA}$, it is the inverse of the present value of an annuity of one dollar a year starting one year later.

The fair rate is individual-specific, and in the model it varies through health. Survival is lower in poor health and with a disability at every age (panel B of \autoref{fig:firststep}), so an individual in worse health typically has lower conditional survival and a higher $\delta^{*}(a)$ (Appendix~\ref{app:derivations} gives a sufficient condition). At a common credit, a one-year delay saves the program more per dollar of PIA when the claimant's conditional survival is lower, and the same delay can save money for one claimant and cost money for another. The fair rate also depends on the length of the horizon over which the higher benefit is paid, which in the model ends at $T=85$. Whether a delay is privately valuable is a separate question: the individual discounts with $\beta$ and her marginal utility of consumption, values the longevity insurance of a larger annuity, and, if she keeps working, avoids the withholding of benefits by the earnings test.

Consider an increase in the DRC from $\delta$ to $\delta+d\delta$, with the reductions for early claiming unchanged. Its direct cost has a mechanical part, the higher benefit paid to inframarginal late claimers, which is nonnegative because the reform raises the benefit at every claiming age after the NRA and lowers it at none, and a behavioral part, the change in the present value of benefits of individuals whose claiming age changes. Each such change is a sequence of one-year steps valued by \eqref{eq:claimnpv} or its analog before the NRA, so the behavioral part is negative when the induced delays occur at ages and for individuals for whom the increment of the new schedule is below its actuarially fair value, $\delta^{*}(a)$ after the NRA (Appendix~\ref{app:derivations}). The direct cost is negative if and only if the behavioral part is negative and exceeds the mechanical cost in absolute value. The indirect cost adds the taxes paid during longer careers and the resulting changes in DI and transfers. The earnings test adds an interaction: benefits withheld after the NRA are credited at the DRC, so a higher DRC lowers the implicit tax on work after the NRA and raises the future benefits of those whose benefits are withheld. The symmetric reforms AL1--AL9 also raise the reduction for early claiming; their mechanical effect is a saving on early claimers and a cost on late claimers, and the same identity values the induced changes in claiming ages.

An increase in the NRA works differently. It shifts the schedule to the right, so that the benefit falls at every claiming age from 62 to 70 (\autoref{tab:schedules}). Its mechanical effect is a cut for every claimant, a saving whatever the claiming ages, to which behavioral responses add or from which they subtract. It does not raise the credit for delay beyond the new NRA; it changes the year-to-year increments of the schedule only by moving the boundary between the early-claiming reductions and the DRC. Its welfare cost falls on all claimants, including those in poor health, whose higher fixed cost of work makes it costlier for them to offset the cut by working longer. Because DI is paid until the NRA without the reduction for early claiming, an NRA increase also lengthens the period over which individuals with a disability can receive DI. A higher DRC rewards delay and raises only the benefits of those who claim after the NRA; a higher NRA mandates delay by lowering the benefits of everyone, including those who do not delay. Section~\ref{sec:results_claiming} shows what this difference implies for welfare and the budget.
